\documentclass[10pt,a4paper]{amsart}
\usepackage[margin=19mm]{geometry}
\usepackage{amsmath,amssymb,mathtools}
\usepackage[scr=rsfs,cal=euler]{mathalfa}
\usepackage{xfrac}
\usepackage[unicode,hidelinks,bookmarks=false]{hyperref}
\newcommand{\ie}{i.e.\ }
\newcommand{\cf}{cf.\ }
\newcommand{\resp}{respectively\ }
\newcommand{\Subsection}{\subsection}
\providecommand{\nobreakdash}{\nobreak\hbox{-}}
\setlength{\emergencystretch}{2em}
\multlinegap0pt
\usepackage{amssymb}
\usepackage{dsfont}
\usepackage{xcolor}
\usepackage{tcolorbox}
\usepackage{tikz}
\newtheorem{theorem}{Theorem}
\newcommand{\AG}{\mathrm{A_G}}
\newcommand{\AR}{\mathrm{A_R}}
\newcommand{\ICG}{\mathrm{IC_G}}
\newcommand{\ICR}{\mathrm{IC_R}}
\newcommand{\eqdef }{\ensuremath{\stackrel{\mbox{\upshape\tiny def.}}{=}}}
\title{Adaptivity Under Realizability Constraints:\\ Comparing In-Context and Agentic Learning}
\author{Anastasis Kratsios, A. Martina Neuman, Philipp Petersen}
\date{Source: arXiv:2605.04995v1 (CC BY 4.0)}
\begin{document}
\maketitle

\noindent\emph{Source and licence.} Adaptivity Under Realizability Constraints:\\ Comparing In-Context and Agentic Learning. Version arXiv:2605.04995v1 (CC BY 4.0), distributed by arXiv under the Creative Commons Attribution 4.0 International licence, http://creativecommons.org/licenses/by/4.0/, as recorded in the arXiv licence metadata for this exact version and shown on the abstract page https://arxiv.org/abs/2605.04995v1. Full licence notice: \texttt{licenses/arxiv-2605.04995-CC-BY-4.0.txt}.\par\smallskip

\noindent\textit{Editorial scope.} Counted unit: the article's theorem thm:address (source0.tex lines 1079--1097). The article's complete original proof of it is delivered with it (source0.tex lines 1647--1800, one \texttt{proof} environment of 154 lines, which the article places away from the statement). No mathematics is written, repaired or re-derived: the statement and the proof are the article's own lines, in the article's own notation, and no retained line is substituted or re-flowed.  Retained uncounted as the source-local context the proof needs: lines 1055--1057; lines 1069--1075.  Nothing was omitted from any retained span, so the package carries no omission record.  No retained line contains a citation, so the wrapper carries no bibliography and no bibliography entry.  No retained line contains a figure environment, an image-inclusion command or drawing code, so no image is delivered and the package declares no asset dependency.  Editorial formatting only, in addition: an AMS \texttt{amsart} preamble that restates the article's own declarations and macros used by the retained lines, namely the environments \texttt{theorem} on separate counters, together with the standard packages those lines need.  The retained environments print in this document as Theorem 1 in order of appearance, whereas the article prints the same items with the numbering of its own full document.  The complete native source of the article as distributed in the arXiv e-print for this exact version is bundled unchanged as \texttt{sources/199914/source0.tex}.
\begin{equation} \label{hats}
    \widehat s\eqdef  (\tau(s_1),\dots,\tau(s_{N-1})).
\end{equation}

\begin{equation} \label{eq:addrtasks}
    f^{\mathrm{addr}}_{s,\beta}(x)
    \eqdef 
    \sum_{i=1}^{N-1} s_i h_{q_i}(x)
    +
    \beta h_{q^\ast(s)}(x),
\end{equation}

\begin{theorem}
\label{thm:address}
Let $m,N\in\mathbb{N}$ with $N\geq 3$.
For the family $\mathcal{T}_{N,m}^{\mathrm{addr}}$, the following hold.
\begin{enumerate}
    \item[\rm (i)] %There exists a general agentic learner that first reads the static values at
    %$q_1,\dots,q_{N-1}$ and then makes one additional adaptive query, reconstructing every task exactly.
    There exists a general agentic learner with a query budget $N$ consisting of $N-1$ fixed queries followed by one adaptive query, that reconstructs every task exactly.
    \item[\rm (ii)] Every in-context learner with the sample budget $N$ incurs worst-case
    $L^\infty([0,1])$ approximation error at least $1/2$.
    \item[(iii)] Every ReLU-realizable agentic learner with query budget \(N\) and query maps implemented by ReLU networks with at most $m$ weights incurs worst-case 
    \(L^\infty([0,1])\) approximation error at least \(1/2\). 
\end{enumerate}
Consequently, 
\begin{equation*}
    \ICG<\AG \quad\text{ and }\quad \ICR=\AR.
\end{equation*}
%Part~\textup{(i)} together with part~\textup{(ii)} yields $\ICG<\AG$. Part~\textup{(iii)} $\ICR = \AR$ since the worst-case error 1/2 can be achieved by a constant approximator.
\end{theorem}

\begin{proof}[Proof of Theorem~\ref{thm:address}]
We prove the three assertions in turn.

\medskip
\noindent
\emph{Proof of (i).}
We construct a general agentic learner $\Psi$ as follows.
For a task $f^{\mathrm{addr}}_{s,\beta}$, the %The general agentic 
learner first queries the fixed points $q_1,\dots,q_{N-1}$. By construction,
this reveals the vector
\[
\bigl(f^{\mathrm{addr}}_{s,\beta}(q_1),\dots,f^{\mathrm{addr}}_{s,\beta}(q_{N-1})\bigr)=s.
\]
The learner then computes $\widehat s$ from \eqref{hats} and queries $g_m$ at $\widehat s$ to obtain
$
    g_m(\widehat s)=q^\ast(s)
$.
Next, the learner makes one additional query of $f^{\mathrm{addr}}_{s,\beta}$ at $q^\ast(s)$. 
Since the supports of the static hats $h_{q_i}$ lie in $[0,1/2+\delta]$
while $q^\ast(s)\in[2/3,1]$, we get from \eqref{eq:addrtasks} a corresponding response
%the static hats vanish there, so
\[
    f^{\mathrm{addr}}_{s,\beta}(q^\ast(s))=\beta.
\]
The learner $\Psi$ now knows both $s$ and $\beta$, and therefore reconstructs the task exactly by
outputting
\[
%\widehat f(x)
\Psi(f^{\mathrm{addr}}_{s,\beta})(x)
=\sum_{i=1}^{N-1} s_i h_{q_i}(x)+\beta h_{q^\ast(s)}(x) = f^{\mathrm{addr}}_{s,\beta}(x).
\]

\medskip
\noindent
\emph{Proof of (ii).}
Let $\Psi$ be an in-context learner with fixed sample points $\xi_1,\dots,\xi_N\in[0,1]$.
Because $\delta<1/(6N)$, the union of the intervals $[\xi_i-\delta,\xi_i+\delta]$ has total length
strictly smaller than $1/3$, so there exists a point $y\in[2/3,1]$ with
$|y-\xi_i|>\delta$ for every $i=1,\dots,N$.
Since $\tau$ is surjective onto $[0,2]$, and $g_m$ is surjective onto $[2/3,1]$, we may choose $s\in[0,1]^{N-1}$ such that $y=g_m(\widehat s)=q^\ast(s)$.

Now consider the two tasks $f^{\mathrm{addr}}_{s,0}$ and $f^{\mathrm{addr}}_{s,1}$. Because the moving hat $h_y=h_{q^\ast(s)}$
%spike 
is supported inside $[y-\delta,y+\delta]$,
both tasks have identical sample values at every
$\xi_i$, yielding the same context 
\begin{equation*}
    C_N(f^{\mathrm{addr}}_{s,0};\xi_1,\dots,\xi_N)
    = C_N(f^{\mathrm{addr}}_{s,1};\xi_1,\dots,\xi_N).
\end{equation*}
As the in-context learner $\Psi$ depends only on this context, it produces the same approximation for both tasks: $\Psi(f^{\mathrm{addr}}_{s,0})=\Psi(f^{\mathrm{addr}}_{s,1})$. 
On the other hand,
\[
f^{\mathrm{addr}}_{s,1}(x)-f^{\mathrm{addr}}_{s,0}(x)=h_{q^\ast(s)}(x),
\]
so at $x=q^\ast(s)$ the two tasks differ by $1$. 
Hence, by the triangle inequality, at least one of the approximations $\|\Psi(f^{\mathrm{addr}}_{s,i})-f^{\mathrm{addr}}_{s,i}\|_{L^{\infty}([0,1])}$ has error
%Therefore one of the two approximation errors is
at least $1/2$.
Since $\Psi$ is arbitrary, we conclude that every in-context learner with sample budget $N$ incurs worst-case $L^{\infty}([0,1])$ approximation error at least $1/2$ on $\mathcal T_{N,m}^{\mathrm{addr}}$.

\medskip
\noindent

\emph{Proof of (iii).}
We first note that the %modified 
address map
$
s\mapsto q^\ast(s)=g_m(\widehat s)
$
is %still 
hard for ReLU networks with at most \(m\) weights. Indeed, suppose that a ReLU network
\(\Phi:[0,1]^{N-1}\to[0,1]\) with at most \(m\) weights satisfies
\[
\|\Phi-g_m\circ \widehat{(\cdot)}\|_{L^\infty([0,1]^{N-1})}<\frac16.
\]
On the subcube \([0,1/2]^{N-1}\), we have \(\tau(t)=2t\), and therefore
\[
    q^\ast(s)=g_m(\widehat{s})=g_m(2s_1,\dots,2s_{N-1}).
\]
Hence the ReLU network
\[
\widetilde \Phi(u_1,\dots,u_{N-1})
\eqdef 
\Phi\left(\frac{u_1}{2},\dots,\frac{u_{N-1}}{2}\right)
\]
has at most the same number of weights and satisfies
\[
\|\widetilde \Phi-g_m\|_{L^\infty([0,1]^{N-1})}<\frac16,
\]
contradicting the choice of \(g_m\).

Now let \(I_i\eqdef  \operatorname{supp}h_{q_i}\), for \(i=1,\dots,N-1\), and
$
    I_*(s)\eqdef  \operatorname{supp}h_{q^\ast(s)}
$.
Assume, for contradiction, that an \(N\)-query ReLU-realizable agentic learner with query maps implementable by ReLU networks with at most \(m\) weights achieves worst-case $L^{\infty}([0,1])$ approximation error strictly smaller than \(1/2\) on $\mathcal T_{N,m}^{\mathrm{addr}}$.

We first claim that the learner must query points in every static support \(I_i\). 
If not, there exists some \(i\) such that no query point falls into \(I_i\). Choose two \(s,s'\in[0,1]^{N-1}\) which agree in all coordinates except the \(i\)th one and satisfy
\[
s_i\neq s_i',
\qquad
\tau(s_i)=\tau(s_i').
\]
For instance, take \(s_i=0\) and \(s_i'=1\). Then $\widehat s=\widehat{s'}$ and therefore
\[
q^\ast(s)=g_m(\widehat s)= g_m(\widehat{s'})=q^\ast(s').
\]
Taking the same value of \(\beta\), the two tasks
\(f^{\mathrm{addr}}_{s,\beta}\) and \(f^{\mathrm{addr}}_{s',\beta}\)
have the same moving spike at $q^\ast(s)=q^\ast(s')$ and the same static spikes at $q_j$, for $j\neq i$.
Thus, they differ \emph{only} on the static support \(I_i\).
On the one hand, since the learner \emph{never} queries in \(I_i\), it receives the same %transcript on 
context from both tasks and hence produces the same approximation. 
On the other hand, the two tasks \(f^{\mathrm{addr}}_{s,\beta}\) and \(f^{\mathrm{addr}}_{s',\beta}\) differ by \(1\) at \(q_i\), by construction.
It follows that at least one of the corresponding approximation errors incurred by the learner is at least \(1/2\), a contradiction.

We next claim that the learner must query points in the moving support \(I_*(s)\). Otherwise, the two tasks \(f^{\mathrm{addr}}_{s,0}\) and \(f^{\mathrm{addr}}_{s,1}\) generate identical 
contexts to the learner, by the now routine argument. %transcripts, because they differ only on \(I_*(s)\). 
The learner would therefore output the same approximation for both tasks. However, since
\[
f^{\mathrm{addr}}_{s,1}(q^\ast(s))
-
f^{\mathrm{addr}}_{s,0}(q^\ast(s))
=
1,
\]
we again conclude that one of the two approximation errors is at least \(1/2\), a contradiction.

Thus, for every \(s\), the \(N\) queries are fully accounted for: \(N-1\) queries go to the static supports and one query goes to the moving support. Since the static supports are fixed, known in advance, pairwise disjoint, and carry independent coefficients, sampling them in arbitrary order reveals exactly the corresponding coordinates \(s_i\).
Therefore, for the lower bound, we may reorder the query points and assume that the learner first reads the static values \(s_1,\dots,s_{N-1}\) and then queries the induced moving support. %choosing its unique moving-support query.
Consequently, the last %moving 
query map is given by a ReLU neural network
\[
\Phi:[0,1]^{N-1}\to[0,1]
\]
with at most \(m\) weights. Since this querying must hit $I_*(s)=\operatorname{supp}h_{q^\ast(s)} = [q^\ast(s)-\delta, q^\ast(s)+\delta]$
for every \(s\), we have
\[
|\Phi(s)-q^\ast(s)|\leq \delta
\quad
\text{ for all }\quad s\in[0,1]^{N-1}.
\]
Therefore
\[
\|\Phi-g_m\circ \widehat{(\cdot)}\|_{L^\infty([0,1]^{N-1})}
\leq \delta.
\]
%Choosing \(\delta<1/6\) contradicts the hardness of the modified address map. 
Since $\delta<1/(6N)\leq 1/18$, this contradicts the hardness of the address map $s\mapsto q^\ast(s)=g_m(\widehat s)$, as established at the beginning of the proof. Hence, every ReLU-realizable agentic learner of this form incurs worst-case \(L^\infty([0,1])\) approximation error at least \(1/2\) on $\mathcal T_{N,m}^{\mathrm{addr}}$.

Finally, parts~(i),~(ii) together yield $\ICG<\AG$, while part~(iii) implies $\ICR = \AR$ since a learner with constant predictor $x\mapsto 1/2$ achieves worst-case error $1/2$.
\end{proof}

\end{document}
